\chapter{Numerical Experiments}




In this chapter, we will present some numerical experiments and subsequent comparative
analysis of DOD-based and POD-based ROMs, as well as the stationary informed CoLoRA architecture
for which we will only produce some qualitative results in a later subsection. The structure of this chapter
will be the following: 

First, we will provide the reader with some specifications of the 
gathering of data snapshots and the general training 
workflow employed for the statistical learning of the neural network weights for the
corresponding models. 
This is uniform in context of the different ROMs to retain comparability. As a next step, we
want to introduce two examples, a one dimensional heat advection problem with geometric initial condition
and a advection-diffusion problem with constant Neumann condition. In both of these, we will
graph the KnW as given by Definition (\ref{def: linear kolmogorov}), or at least a numerical approximation of
it, then show one qualitative realization of a neural network architecture with trained weights for the
POD-DL-ROM, the DOD+DFNN and the DOD-DL-ROM plottet against a FOM solution, and finally
set up a comparison in relative error in relation to the active weights to compare the theoretical result
to numerical ones.

In order to connect to the brief discussion of Statical Learning in Subsection (\ref{subsec: statistical learning}),
we will start with the specific definitions of the training procedure, after attaining a set of data in
the form of
\begin{equation*}
    \left\{ (\mu_i, \nu_j, k \Delta t), 
    u_h^{\mu_i, \nu_j, k \Delta t}\right\}_{(i,j,k) = (1, 1, 1)}^{(N_{s_1}, N_{s_2}, N_t)} 
    \subset (\Pcal_1 \times \Pcal_2 \times [0,T]) \times \R^{N_h}.
\end{equation*}
Here, we generate the solutions using the \texttt{pymor} module \cite{pymor2016}, and its builtin discretizers relying
on \texttt{gmsh} meshing for the second, more complex, example. In general we use a CG discretizer technique
and its native solver, but this can of course be modified freely. Crucially, 
we perform the POD and each of the error-estimates employing the gram matrix generated by
the $W^{1, 2}_0$-norm on the deaffinized solution space. For this, we subtract a 
Dirichlet shift before training and re-add it afterwards, so that any model is being trained
on the correct space. Hence any difference between solutions in $\norm{\cdot}$ should be seen
as close in a $W^{1, 2}$-sense, since the matrix gram matrix $\G$ is taken from that
inherent inner product. 

The data can be understood 
in terms of solution snapshot matrices $S \in \R^{N_h \times N_{s_1}N_{s_2}N_t}$ and their respective parameter stacks 
$M \in \R^{p \times N_{s_1}}, N \in
\R^{q \times N_{s_2}}$:
\begin{equation*}
    S = \begin{bmatrix}[c|c|c|c|c|c|c]
        u_h^{\mu_1, \nu_1, \Delta t} & \cdots & u_h^{\mu_{N_{s_1}}, \nu_1, \Delta t} & u_h^{\mu_2, \nu_1, \Delta t}
        & \cdots & \cdots & u_h^{\mu_{N_{s_1}}, \nu_{N_{s_2}}, N_t \Delta t}.
    \end{bmatrix}
\end{equation*}
After this, we use a common validation strategy, where we split each batch into a training data
and a validation data, such that $M = [M^\text{train}, M^\text{val}], 
N = [N^\text{train}, N^\text{val}]$ and $S = [S^\text{train}, S^\text{val}]$ and a splitting factor of $\alpha=0.8$.

For the normalization, we use the aforementioned min-max normalization in a generalized sense; for
each parameter set we define 
\begin{equation*}
    M_\text{max}^{i} = \max\limits_{j = 1, \dots, N_{s_1}} M_{ij}^\text{train}, \quad 
    M_\text{min}^{i} = \min\limits_{j = 1, \dots, N_{s_1}} M_{ij}^\text{train}
\end{equation*}
and its subsequent min-max normalization via
\begin{equation} \label{eq: min-max normalization}
    M_{ij}^\text{train} \mapsto \frac{M_{ij}^\text{train} - M_\text{max}^i}{M_\text{max}^i - M_\text{min}^i},
    \quad \text{for } i = 1, \dots, p, \, j = 1, \dots, N_{s_1}.
\end{equation}
For $N^\text{train}$ the analogous approach can be used. The solution batch is normalized in the
same way, after defining the minimum and maximum as
\begin{equation*}
    S_\text{max} = \max\limits_{i = 1, \dots, N_h} \max\limits_{j = 1, \dots, N_s N_t} S_{ij}^\text{train},
    \quad S_\text{min} = \min\limits_{i = 1, \dots, N_h} \min\limits_{j = 1, \dots, N_s N_t} S_{ij}^\text{train}.
\end{equation*}

All offline training phases will be performed with the AdamW optimizer and the CosineAnnealingWarmRestarts schedular 
to achieve adaptive learning rates. Moreover, an early stopping criterion will be imposed upon the
validation loss. The relevant hyperparameters will be discussed for each example seperately,
as we choose them with regard to the specifics of the problem at hand. However, in the following
Algorithm (\ref{alg: dod-dl-rom training})
we describe the exact offline procedure for the DOD-DL-ROM training collectively. The non-DOD part of
the DOD+DFNN is a less complex adaptation of a straight-forward concatenation of some linear layers, 
and can hence be infered from the DOD-DL-ROM learning procedure and its definition in Subsection
(\ref{subsec: dod-based roms}), since the 
a priori treatment of the data is the same. 
The POD-DL-ROM offline and online procedure can be found in \cite{fresca2022}.


For evaluation, we rely on the following discretized version of the relative error $\Ecal_\text{R} \in [0, \infty)$
over a test sample $\{(\mu_i, \nu_i), u_h^{\mu_i, \nu_i}\}_{i=1}^{N_\text{test}}$ with time-trajectories
in $\Tcal$ defined as
\begin{equation*}
    \Ecal_\text{R} = \frac{1}{N_\text{test}} \sum_{i=1}^{N_\text{test}} \left(
        \frac{
            \sqrt{\frac{1}{N_t} \sum_{k=1}^{N_t} \norm{u_h^{\mu_i, \nu_i, k \Delta t} - 
            \hat{u}_h^{\mu_i, \nu_i, k \Delta t}}^2}
        }{
            \sqrt{\frac{1}{N_t} \sum_{k=1}^{N_t} \norm{u_h^{\mu_i, \nu_i, k \Delta t}}^2}
        }
        \right)
\end{equation*}
as well as the pointwise time-trajectory 
error $\epsilon_\text{rel} \in [0, \infty)^{N_h \times N_t}$ for visualization, given by
\begin{equation*}
    \epsilon_\text{rel} = \frac{
        \euknorm{u_h^{\mu_i, \nu_i, k \Delta t} - 
        \hat{u}_h^{\mu_i, \nu_i, k \Delta t}}_\text{pw}
    }{
        \sqrt{\frac{1}{N_t} \sum_{k=1}^{N_t} \norm{u_h^{\mu_i, \nu_i, k \Delta t}}^2}
    }
\end{equation*}
where we set $\euknorm{\cdot}_\text{pw}$ as the pointwise absolute difference of the $N_h$-dimensional 
vector, and $\hat{u}_h^{\mu, \nu, t} \in \R^{N_h}$ denoting some solution prediction. 


Model training is performed using the Python library \texttt{PyTorch}~\cite{pytorch2019} on
a Nvidia Quadro RTX 6000. 
All source code required to reproduce the presented results, accompanied by a concise 
\texttt{README.md}, is openly available at 
\url{https://github.com/dawid-kotowski/doddlrom}.


\begin{algorithm}[H]
    \caption{DOD training algorithm (offline)}
    \label{alg: inner-dod training}
    \begin{algorithmic}[1]
    \Require Parameter matrices $M \in \R^{p\times N_{s_1}}$,
    snapshot matrix $S \in \R{N_h \times N_t N_s}$, 
    training-validation splitting fraction $\alpha$, starting learning rate $\eta^0$, 
    batch size $N_b$, batch number $N_{nb}$, maximum number of epochs $N_{\text{epochs}}$, dimensions $n, N', N_A$
    \Ensure Optimal model parameters $\theta_\text{DOD}^*$
    \State Compute pre-reduction POD basis matrix $\A \in \R^{N_h \times N_A}$ via Definition
    (\ref{def: POD})
    \State Randomly shuffle $M$ and $S$
    \State Split data in $M = [M^{\text{train}}, M^{\text{val}}]$ and 
    $S = [S^{\text{train}}, S^{\text{val}}]$ according to $\alpha$
    \State $S_{N_A}^\text{train} \gets \A^T \G S^\text{train}$ and $S_{N_A}^\text{val} \gets \A^T \G S^\text{val}$
    \State Normalize data in $M$ and $S_{N_A} = [S_{N_A}^{\text{train}}, S_{N_A}^{\text{val}}]$
    \State Initialize inner DOD module $\Phi_{\tilde{\V}}(\theta^0_\text{DOD}) \gets \text{RandomInit}(\theta^0_\text{DOD})$
    \State $n_e = 0$
    \While{(¬early-stopping and $n_e \leq N_{\text{epochs}}$)}
      \For{$k = 1:N_{nb}$}
        \State Sample a batch $(M^{\text{batch}}, N^{\text{batch}}, S_{N_A}^{\text{batch}}) \subset 
        (M^{\text{train}}, N^{\text{train}}, S_{N_A}^{\text{train}})$
        \For{$t = 1:N_t$}
            \State $\tilde{\V}_{M^\text{batch}, t} \gets 
            R_\rho(\Phi_{\tilde{\V}}(\theta^{n_e N_b + k}_\text{DOD}))(M^\text{batch}, t)
            \in \R^{N_b \times N_h \times N'}$
        \EndFor
        \State Accumulate loss $\Lcal_\text{DOD}$ of Definition
        (\ref{def: DOD}) on $(M^{\text{batch}}, S_{N_A}^{\text{batch}})$ 
        and compute $\nabla_\theta \Lcal$
        \State $\theta^{n_e N_{\text{nb}}+k+1} \gets 
        \text{AdamW}(\eta^{n_e}, \nabla_\theta \Lcal, \theta_\text{DOD}^{n_e N_{\text{nb}}+k})$
        \State $\eta^{n_e + 1} \gets \text{CAWR}(\Lcal_\text{DOD})$
      \EndFor
      \State Repeat instructions 9--17 on $(M^{\text{val}}, S_{N_A}^{\text{val}})$ 
      with the updated weights $\theta_\text{DOD}^{n_e N_{\text{nb}}+k+1}$
      \State Accumulate loss $\Lcal_\text{DOD}$ of Definition
        (\ref{def: DOD}) on $(M^{\text{val}}, S_{N_A}^{\text{val}})$ 
      to evaluate early-stopping criterion
      \State $n_e = n_e + 1$
    \EndWhile
    \end{algorithmic}
    \end{algorithm}




    \begin{algorithm}[H]
        \caption{DOD-DL-ROM training algorithm (offline)}
        \label{alg: dod-dl-rom training}
        \begin{algorithmic}[1]
        \Require Parameter matrices $M \in \R^{p\times N_{s_1}}, N \in \R^{q \times N_{s_2}}$, 
        snapshot matrix $S \in \R{N_h \times N_t N_s}$, 
        training-validation splitting fraction $\alpha$, starting learning rate $\eta^0$, 
        batch size $N_b$, batch number $N_{nb}$, maximum number of epochs $N_{\text{epochs}}$, dimensions $n, N', N_A$
        \Ensure Optimal model parameters $\theta_\text{DOD-DL-ROM}^* = (\theta_\text{DOD}^*, \theta_\text{D}^*,
        \theta_\text{DF}^*)$
        \State Compute pre-reduction POD basis matrix $\A \in \R^{N_h \times N_A}$ via Definition
        (\ref{def: POD})
        \State Randomly shuffle $M$ and $S$
        \State Split data in $M = [M^{\text{train}}, M^{\text{val}}], N=[N^\text{train}, N^\text{val}]$ and 
        $S = [S^{\text{train}}, S^{\text{val}}]$ according to $\alpha$
        \State $S_{N_A}^\text{train} \gets \A^T \G S^\text{train}$ and $S_{N_A}^\text{val} \gets \A^T \G S^\text{val}$
        \State Normalize data in $M, N$ and $S_{N_A} = [S_{N_A}^{\text{train}}, S_{N_A}^{\text{val}}]$
        \State Train inner DOD module $\Phi_{\tilde{\V}}(\theta^*_\text{DOD})$ with Algorithm 
        (\ref{alg: inner-dod training})
        \State Initialize DL-ROM weights $(\theta^0_\text{D}, \theta^0_\text{E}, \theta^0_\text{DF}) 
        \gets \text{RandomInit}(\theta^0_\text{D}, \theta^0_\text{E}, \theta^0_\text{DF})$
        \State $n_e = 0$
        \While{(¬early-stopping and $n_e \leq N_{\text{epochs}}$)}
        \For{$k = 1:N_{nb}$}
            \State Sample a batch $(M^{\text{batch}}, N^{\text{batch}}, S_{N_A}^{\text{batch}}) \subset 
            (M^{\text{train}}, N^{\text{train}}, S_{N_A}^{\text{train}})$
            \For{$t =1:N_t$}
                \State $S_{N'}^{\text{batch}, t} \gets \tilde{\V}^T_{M^\text{batch}, t} \cdot 
                S_{N_A}^{\text{batch}, t} \in \R^{N_b \times N'}$
                \State $S_{N'}^{\text{batch}, t} \gets 
                \text{reshape}(S_{N'}^{\text{batch}, t}) \in 
                \R^{N_b \times \sqrt{N'} \times \sqrt{N'}}$
                \State $z^{(n_e N_{\text{nb}}+k), t} \gets
                R_\rho(\Psi_n'(\theta_\text{E}^{n_e N_{\text{nb}}+k}))(S_{N'}^{\text{batch}, t})$
                \State $\hat S_{N'}^{\text{batch}, t} \gets
                R_\rho(\Psi_{N'}(\theta_\text{D}^{n_e N_{\text{nb}}+k}) \odot \Phi_n(\theta_\text{DF}^{n_e N_{\text{nb}}+k}))(M^{\text{batch}}, N^\text{batch}, t)$
                \State $\hat S_{N'}^{\text{batch}, t} \gets
                \text{reshape}(\hat S_{N'}^{\text{batch}, t}) 
                \in \R^{N_b \times {N'}}$
            \EndFor
            \State Accumulate loss $\Lcal_\text{DOD-DL-ROM}$
            from Subsection (\ref{subsec: dod-dl-rom}) on $(M^{\text{batch}}, N^\text{batch}, S_{N_A}^{\text{batch}})$ 
            \State $\theta^{n_e N_{\text{nb}}+k+1} = 
            \text{AdamW}(\eta^0, \nabla_\theta \mathcal{J}, \theta^{n_e N_{\text{nb}}+k})$
            \State $\eta^{n_e + 1} \gets \text{CAWR}(\Lcal_\text{DOD})$
        \EndFor
        \State Repeat instructions 9--24 on $(M^{\text{val}}, N^\text{val}, S_{N_A}^{\text{val}})$ 
        with the updated weights $\theta^{n_e N_{\text{nb}}+k+1}$
        \State Accumulate loss $\Lcal_\text{DOD-DL-ROM}$
            from Subsection (\ref{subsec: dod-dl-rom}) on $(M^{\text{val}}, N^\text{val}, S_{N_A}^{\text{val}})$
        \State $n_e = n_e + 1$
        \EndWhile
        \end{algorithmic}
        \end{algorithm}




\section{Two Visual Presentations}

In the following, we want to show the reader two advection-reaction-diffusion equations, which
meet the criterion of Assumption (\ref{assumption KnW decay}) in a sense;
both have significantly faster KnW decay for fixed parameter $(\mu, t) \in \Theta \times [0, T]$. 


\subsection{Example 1: Wind from a Hole in a complex domain}



We want to start with a simple problem on a relatively complex domain; a rotating wind 
from a hole. To circumvent the CFL-barrier, we impose the time domain to $T=0.008$.
Consider a polygonal domain $\Omega \subset (0,1)^2$
with two circular holes (see Figure \ref{fig:polygon-domain}). 
Let $u^{\mu,\nu}$ be called a strong solution to the problem, if for each 
$\mu \in (\pi+0.5, 2\pi)$ and $\nu \in (0.01,0.08)$ we have

\begin{equation*}
    \begin{cases}
    \displaystyle
    \partial_t u^{\mu,\nu}(x,t)
    + b(\mu)\cdot\nabla u^{\mu,\nu}(x,t)
    - \nabla\!\cdot\!\bigl(\kappa(x;\nu)\nabla u^{\mu,\nu}(x,t)\bigr)
    + \alpha\,u^{\mu,\nu}(x,t)
    = 0, 
    & (x,t)\in \Omega\times(0,T], \\[0.8em]
    \displaystyle
    u^{\mu,\nu}(x,t) = 3, 
    & (x,t)\in \Gamma_D\times(0,T], \\[0.8em]
    \displaystyle
    (\kappa(x;\nu)\nabla u^{\mu,\nu})(x,t)\cdot n(x) = 0, 
    & (x,t)\in \Gamma_N\times(0,T], \\[0.8em]
    \displaystyle
    u^{\mu,\nu}(x,0) = 0, 
    & x\in \Omega,
    \end{cases}
\end{equation*}
where the advection, diffusion and reaction terms are given by
\begin{equation*}
    b(\mu) = 
    \begin{pmatrix}
    40\cos(\mu) \\[0.3em]
    40\sin(\mu)
    \end{pmatrix},
    \qquad
    \kappa(x;\nu) = \nu\,(1-x_1) + x_1,
    \qquad
    \alpha = 10^{-1}.
\end{equation*}
Here, $\Gamma_D$ denotes the inner boundary of the small circular hole centered at $(0.65,0.35)$
with radius $0.05$, on which Dirichlet boundary conditions are imposed,
while $\Gamma_N$ denotes the outer boundary and the larger hole centered at $(0.35,0.65)$
with radius $0.10$, on which homogeneous Neumann boundary conditions are imposed.



  \begin{figure}[h]
    \centering
    

    \begin{tikzpicture}[x=5cm,y=5cm,>=Latex,
        neumann/.style   ={very thick},
        dirichlet/.style ={very thick,dashed},
        coord/.style     ={black,fill=black,circle,inner sep=0.6pt}
      ]
      
      \coordinate (p1)  at (0.00,0.00);
      \coordinate (p2)  at (1.00,0.00);
      \coordinate (p3)  at (1.00,0.20);
      \coordinate (p4)  at (0.80,0.20);
      \coordinate (p5)  at (0.80,0.40);
      \coordinate (p6)  at (1.00,0.40);
      \coordinate (p7)  at (1.00,1.00);
      \coordinate (p8)  at (0.00,1.00);
      \coordinate (p9)  at (0.00,0.20);
      \coordinate (p10) at (0.30,0.20);
      \coordinate (p11) at (0.30,0.10);
      \coordinate (p12) at (0.00,0.10);
      
      \def\hcA{(0.65,0.35)}
      \def\hrA{0.05}   
      \def\hcB{(0.35,0.65)} 
      \def\hrB{0.10}       
      
      \fill[gray!12]
        (p1)--(p2)--(p3)--(p4)--(p5)--(p6)--(p7)--(p8)--(p9)--(p10)--(p11)--(p12)--cycle;
      \fill[white] \hcA circle[radius=\hrA];
      \fill[white] \hcB circle[radius=\hrB];
      
      \draw[neumann]
        (p1)--(p2)--(p3)--(p4)--(p5)--(p6)--(p7)--(p8)--(p9)--(p10)--(p11)--(p12)--cycle;
      
      \draw[dirichlet] \hcA circle[radius=\hrA];
      \draw[neumann]   \hcB circle[radius=\hrB];
      
      \foreach \k in {1,...,12}{
        \path (p\k) node[coord]{};
      }
      
      \begin{scope}[xshift=0.02cm,yshift=0.02cm]
        \draw[neumann] (0.03,0.95) -- ++(0.08,0) node[right,black,scale=0.9]{Neumann};
        \draw[dirichlet] (0.03,0.88) -- ++(0.08,0) node[right,black,scale=0.9]{Dirichlet};
      \end{scope}
      
      \node[scale=0.9] at ($(0.65,0.35)+(0, -0.1)$) {$\Gamma_D$};
      \node[scale=0.9] at ($(1,0.65)+(-0.07, 0)$) {$\Gamma_N$};
      \node[scale=0.9] at ($(0,0.35)+(+0.07, 0)$) {$\Gamma_N$};
      \node[scale=0.9] at ($(0.65,1)+(0, -0.07)$) {$\Gamma_N$};
      \node[scale=0.9] at ($(0.65,0)+(0, +0.05)$) {$\Gamma_N$};
      \node[scale=0.9] at ($(0.35,0.65)+(0,  -0.15)$) {$\Gamma_N$};
    \end{tikzpicture}


    \caption{Polygonal subdomain of $(0,1)^2$ with Dirichlet and Neumann boundaries.}
    \label{fig:polygon-domain}
  \end{figure}


\subsection{Example 2: Wind Generation in a plain domain}




For the second example, we will look at a much simpler square domain $\Omega = (0, 1)^2$ 
with last time point $T = 1$. Let $u^{\mu, \nu}$ be called a strong solution to the
problem, if for each $\mu = (\mu_1, \mu_2, \mu_3) \in (0.3, 0.7)^3$ and $\nu \in (0.002, 0.005)$ we have 

\begin{equation*}
    \begin{cases}
    \displaystyle
    \partial_t u^{\mu,\nu}(x,t)
    + b(\mu_3)\cdot\nabla u^{\mu,\nu}(x,t)
    - \nu\,\Delta u^{\mu,\nu}(x,t)
    = f(x;\mu), 
    & (x,t)\in \Omega\times(0,T], \\[0.8em]
    \displaystyle
    (\nu\,\nabla u^{\mu,\nu})(x,t)\cdot n(x) = 0, 
    & (x,t)\in \partial\Omega\times(0,T], \\[0.8em]
    \displaystyle
    u^{\mu,\nu}(x,0) = 0, 
    & x\in \Omega,
    \end{cases}
\end{equation*}
where the advection and right-hand side terms are given by
\begin{equation*}
    b(\mu_3) =
    \begin{pmatrix}
    \cos\left(\tfrac{\pi}{100\,\mu_3}\right) \\[0.3em]
    \sin\left(\tfrac{\pi}{100\,\mu_3}\right)
    \end{pmatrix},
    \qquad
    f(x;\mu) = 10 \exp\left(-\frac{(x_1-\mu_1)^2+(x_2-\mu_2)^2}{0.07^2}\right).
\end{equation*}

